
%%%%%%%%%%%%%%%%%%%%%%%%%%%%%%%%%%%%%%%%%%%%%%%%%%%%%%%%%%%%%
\section{模型的改进与推广}

\subsection{模型的改进}

针对模型在训练集上过拟合的问题，未来可引入L1正则化（Lasso）或L2正则化（Ridge）对GBR的损失函数进行约束，同时采用早停（early stopping）机制在验证集误差不再下降时终止训练。此外，可将当前的单一GBR模型升级为Stacking集成策略，将随机森林、GBR和线性回归的预测结果进行元学习融合，进一步提升预测精度。针对特征获取的局限性，可探索将宏观指标（如地区GDP增速、消费信心指数）和外部数据（如天气、社交媒体热度）纳入特征体系，增强模型的适用场景和鲁棒性。

\subsection{模型的推广}

本文构建的零售销售额预测与分析框架不仅适用于案例中的零售企业，还可推广至餐饮连锁、快消品零售、百货商场等具有类似经营数据结构的行业。在餐饮场景中，只需将客流量替换为翻台率、将客单价替换为人均消费，即可复用本文的分析—预测—策略方法框架。在电商场景中，促销力度和线上占比两个维度已天然适用，仅需将门店经营指标映射为平台运营指标即可。本文提出的多因素归因评估模型和情景分析策略方法具有较强的通用性和跨行业迁移能力，为数据驱动的精细化运营提供了可借鉴的方法论参考。
